\documentclass[14pt,aspectratio=1610]{beamer}
%\documentclass[14pt,aspectratio=1610,handout]{beamer}
\usepackage{csu-theme}
% \usepackage{linked-lists}
\usepackage{binary-trees}
\usepackage{booktabs}

\forestset{
  default preamble={%
    for tree={%
      treeNode,%
      %font=\small,%
      % l sep=0.2em,%
      % level distance=0.5em,%
      % outer sep=0pt,
      % outer ysep=0pt,
      l=2.4em,
    },%
    %
    % Don't show circle around leaves (for null pointers)
    for leaves={%
      justText},
  },%
}

\title{\Huge{} Binary Search Trees\\Generic Traversal}
\subtitle{CSCI 315: \textit{Data Structure Analysis}\\Chapter 19}
%\author{Sean T. Hayes, Ph.D.}
\institute{Charleston Southern University}

%% Comment out date to use default (today)
% \date{March 6, 2025}
\subject{Software Programming}
%\keywords{}
\graphicspath{{../imgs/binary-tree}}

\setlength{\parskip}{0.5em}

\begin{document}

\begin{frame}
  \titlepage{}
\end{frame}

\section{Generic Traversal}
\begin{frame}{Generic Traversal}
  \begin{itemize}[<+->]
    \item
      In a traversal algorithm, ``visiting'' may mean different things.
      \begin{itemize}
        \item For example, outputting the values, updating the values, and calculating results.
      \end{itemize}
    \item
      \emph{Problems}:
      \begin{itemize}
        \item
          Repeatedly implementing a specific traversal algorithm for each function.
        \item 
          Must have access to the internals of the tree.
      \end{itemize}
    \item
      \emph{The Solution}: Write a generic traversal function that accepts a function as a parameter.
  \end{itemize}
\end{frame}

\begin{frame}{Generic Traversal}
  \begin{itemize}[<+->]
    \item
      In C++, a function name without parentheses is considered a pointer to the function.
    \item
      To specify a function as a formal parameter to another function,
      \begin{enumerate}
        \item specify the function type (e.g., \lstinline{void}),
        \item the name as a pointer (e.g., \lstinline{(*visit)}),
        \item the parameter types (e.g., \lstinline{(int, double)}).
      \end{enumerate}
    \item For example,\\
      \lstinline{void funcParamFunc1(void (*visit) (int));}\\[1em]
      \lstinline{template<typename Type>}\\
      \lstinline{void funcParamFunc2(void (*visit) (Type&));}
  \end{itemize}
\end{frame}

\begin{frame}[fragile]{Generic Traversal Example}
A recursive member function that accepts a node to tree to traverse and a function to call at each node.
\begin{lstlisting}[basicstyle=\ttfamily\small]
template<typename Type>
void BinarySearchTree<Type>::inorder(Node *pRoot,
      void (*visit)(Type& value)) {
  if (pRoot != nullptr)  {
    inorder(pRoot->left, visit);

    visit(pRoot->value);

    inorder(pRoot->right, visit);
  }
}
\end{lstlisting}
\end{frame}

\section{Function Objects}
\begin{frame}{Function Objects (a.k.a. Functors)}
  \begin{itemize}[<+->]
    \item
      In short, and object that looks like a function, but that\\
      can maintain state.
    \item
      A function object is an object class/struct that overloads \lstinline{operator()}.
    \item
      Can be passed to a generic traversal function instead of a function pointer.
  \end{itemize}
\end{frame}

\begin{frame}[fragile]{Function Object: Example Definition}
\begin{lstlisting}[basicstyle=\ttfamily\small]
template <typename Type>
struct SumValues {
  // Constructor initializes sum to 0
  SumValues() : sum(0) {}

  // Overloaded operator() to accumulate sum
  void operator()(Type data) {
    sum += data;
  }

  Type sum;
};
\end{lstlisting}
\end{frame}

\begin{frame}[fragile]{Function Object: How it Works}
\begin{lstlisting}[basicstyle=\ttfamily\small]
SumValues funObj;

funObj(10);
funObj(20);

std::cout << funObj.sum << '\n'; // 30
\end{lstlisting}
\end{frame}

\begin{frame}[fragile]{Function Object: Example Generic Traversal}
\begin{lstlisting}[basicstyle=\ttfamily\small]
template<typename Type>
template<typename Callable>
void BinarySearchTree<Type>::preorderTraversal( Node* pCurr, Callable& callable)
{
  if (pCurr != nullptr)
  {
    callable(pCurr->value); // visit the node
    preorderTraversal(pCurr->pLeft, callable);
    preorderTraversal(pCurr->pRight, callable);
  }
}
\end{lstlisting}
\end{frame}

\begin{frame}[fragile]{Function Object: Example Generic Traversal}
\begin{lstlisting}[basicstyle=\ttfamily\small]
template<typename Type>
template<typename Callable>
void BinarySearchTree<Type>::preorderTraversal( Callable& callable)
{
  preorderTraversal(mpRoot, callable);
}
\end{lstlisting}
\end{frame}

\begin{frame}[fragile]{Function Object: Example Generic Traversal}
\begin{lstlisting}[basicstyle=\ttfamily\small]
BinarySearchTree<int> tree{VALUES, SIZE};
SumValues<int> sumObj; // Instantiate function object

tree.preorderTraversal(sumObj);

std::cout << "Sum: " << sumObj.sum << '\n';
\end{lstlisting}
\end{frame}

\section{Self-Balancing}
\begin{frame}{Self-Balancing Binary Search Trees}
\begin{itemize}[<+->]
  \item
    The worst-case performance for searching is \(O(n)\).
  \item
    \emph{Self-Balancing Binary Trees} automatically \\
    minimize the tree's height when performing\\
    insertion and deletion operations.
  \item
    Searching (and therefore insertion and deletion) becomes \(O(\log_2{n})\) (i.e., asymptotically optimal) in the worst case (and may average \(O(1)\) for some operation depending on the implementation).
  \item
    \emph{Disadvantage}: Algorithms are more complicated (especially if duplicates are allowed).
\end{itemize}
\end{frame}

\begin{frame}{Examples of Self-Balancing Binary Search Trees}
\begin{itemize}
  \item
    2--3 --- nodes hold 2 or 3 keys; splits and merges balance.
  \item
    AVL --- strictly height-balanced using rotations when child subtree heights differ by more than 1.
  \item
    Red Black --- guarantees \(O(\log_2{n})\) height with fewer rotations on updates.
  \item 
    AA --- red-black tree with simpler level rules and rotations.
  \item
    Scapegoat --- rebuilds unbalanced subtrees when size/height invariants are violated.
  \item
    Splay --- moves accessed nodes to the root via rotations.
  \item
    Treap --- BST order with random heap priorities.
  \item
    etc.
\end{itemize}
\end{frame}

\begin{frame}{Example of Balancing using Rotation}
\centering

\begin{forest}
  [,phantom, for tree={fit=band, l sep-=1em, minimum size=1.5em, font=\normalsize}, for leaves={treeNode}, for tree={s sep=0.5cm}, s sep=1.5cm
    [3
      [1, name=p1
        [,justText, no edge, for current={no edge}, name=hid1] {
          \onslide<2->{\draw[<-,dotted, very thick] (hid1) to[bend left] node[midway,left] {rotate} (p1);}
        }
        [2,name=c1] {
          \onslide<2->{\draw[->,dotted, very thick] () to[bend right] node[midway,right] {rotate} (p1);}
        }
      ]
      [, phantom]
    ]
    [3, name=p2, for tree={visible on=<3->}
      [2
        [1]
        [, phantom]
      ]{
        \onslide<4->{\draw[->,dotted, very thick] () to[bend left] node[midway,left] {rotate} (p2);}
      }
      [,justText, for current={no edge}, name=hid2] {
        \onslide<4->{\draw[<-,dotted, very thick] () to[bend right] node[midway,right] {rotate} (p2);}
      }
    ]
    [2, for tree={visible on=<5->}
      [1]
      [3]
    ]
  ]
\end{forest}
\end{frame}

\section*{Lab 15}
\begin{frame}{Lab 15: Binary Trees (Part 2)}
\begin{center}\Large
Let's take a look at the lab based on today's lecture.
\end{center}
\end{frame}
\end{document}
